% old_appendix.tex
% Subsections cut from main.tex but kept around for reference / potential restoration.
% Not \input from main.tex — this file is storage only.
%
% Contents:
%   1. App. B.4 — Robustness across probe-training positions
%      (was \subsection at main.tex:481, label app:cross-training-positions, fig:cross-training)
%   2. App. K.3 — Stated label and executed content can come apart
%      (was \subsection at main.tex:952, label app:eot-flip-quality, fig:eot-flip-quality)


% ============================================================
% 1. App. B.4 — Robustness across probe-training positions
% ============================================================

\subsection{Robustness across probe-training positions}
\label{app:cross-training-positions}

The main-text probe is trained at the end-of-turn token. As a robustness check, we re-train the same probe methodology at two alternative positions from App.~\ref{app:token-selection}: the role-marker token and the task-averaged position. Truth, harm, and democrat-prompt politics discriminate with the same sign across all three training positions (Fig.~\ref{fig:cross-training}). Under the republican prompt the readout moves toward inversion: the role-marker probe shows a soft sign flip, while the end-of-turn and task-averaged probes narrow toward zero without crossing it. All bars use the canonical L\gemmaClassificationProbeLayer{} on Gemma.

\begin{figure}[!t]
  \centering
  \includegraphics[width=0.85\linewidth]{plot_042126_cross_training_selector.png}
  \caption{\textbf{Preference vectors obtained at different token positions also generalise to out-of-distribution preferences.} We train linear probes on different kinds of tokens (end-of-turn, role marker, task-average see App.~\ref{app:token-selection}), otherwise using the same methodology. We then evaluate those probes on the usual end-of-turn token. We find that these probes also generalise to discriminating between true and false statements, and also modulating scores under various system prompts. Positive = first-listed class scores higher (true, harmful, left). All three probes agree in sign on truth, harm, and democrat-prompt politics. Under the republican prompt the readout moves toward inversion: the role-marker probe shows a soft sign flip, while the end-of-turn and task-averaged probes narrow toward zero without crossing it.}
  \label{fig:cross-training}
\end{figure}


% ============================================================
% 2. App. K.3 — Stated label and executed content can come apart
% ============================================================

\subsection{Stated label and executed content can come apart}
\label{app:eot-flip-quality}

A flip in the scoring above is a switch in the model's \emph{stated} label (\texttt{Task A:} / \texttt{Task B:}). It does not necessarily mean the model executed the chosen task's content. To audit, we re-generated 200 randomly sampled flipped completions and used an LLM judge to independently classify the stated label and the executed content (Fig.~\ref{fig:eot-flip-quality}).

Most flips ($\sim 72\%$) are \textbf{genuine}: label and content both flipped. A non-trivial $\sim 20\%$ are \textbf{label-only}: the model says ``Task B:'' and then proceeds to write Task A's content; a small remainder are refusals or ambiguous. The EOT token primarily controls the \emph{routing} decision (which label is emitted first), and content execution can fall back on information elsewhere in the prompt.

The dissociation is asymmetric on harm. On one-harmful pairs, patching \emph{toward} the harmful task yields mostly genuine flips ($\sim 84\%$); patching \emph{away} from the harmful task yields mostly label-only flips ($\sim 51\%$ label-only, $\sim 38\%$ genuine, the rest refusals). Harmful task content is stickier than benign content under EOT-only override. Harm-relevant features in the prompt task tokens remain causally active downstream even when the routing signal is reversed.

\begin{figure}[!t]
  \centering
  \includegraphics[width=0.78\linewidth]{plot_030626_flip_type_by_harm.png}
  \caption{\textbf{Flip type by harm composition of the pair.} Genuine content flips dominate on benign pairs ($83\%$). One-harmful pairs show a large gap between toward-harm and away-from-harm patching: toward-harm flips are mostly genuine, away-from-harm flips are mostly label-only. Pair counts: neither harmful $n=99$, one harmful $n=89$, both harmful $n=12$.}
  \label{fig:eot-flip-quality}
\end{figure}
